% SPEC-TH-001 — Thalamus Orchestrator

% Shared preamble for all spec documents.
% Include from each spec: % Shared preamble for all spec documents.
% Include from each spec: % Shared preamble for all spec documents.
% Include from each spec: \input{../preamble.tex} (adjust depth).

\documentclass[11pt,a4paper]{article}

\usepackage[utf8]{inputenc}
\usepackage[T1]{fontenc}
\usepackage{lmodern}
\usepackage[a4paper,margin=2.3cm]{geometry}
\usepackage{parskip}
\usepackage{microtype}
\usepackage[dvipsnames]{xcolor}
\usepackage{enumitem}
\usepackage{listings}
\usepackage{tcolorbox}
\usepackage{titlesec}
\usepackage{fancyhdr}
\usepackage{array}
\usepackage{longtable}
\usepackage{booktabs}
\usepackage{amsmath}
\usepackage{amssymb}
\usepackage{hyperref}
\usepackage{cleveref}

\tcbuselibrary{skins,breakable}

\definecolor{specBlue}{RGB}{30,64,132}
\definecolor{specGray}{RGB}{90,90,90}
\definecolor{specAccent}{RGB}{198,40,40}
\definecolor{specGreen}{RGB}{30,110,60}

\hypersetup{
  colorlinks=true,
  linkcolor=specBlue,
  urlcolor=specBlue,
  citecolor=specBlue,
  pdfborder={0 0 0}
}

\titleformat{\section}{\Large\bfseries\color{specBlue}}{\thesection}{0.75em}{}
\titleformat{\subsection}{\large\bfseries\color{specBlue!80}}{\thesubsection}{0.6em}{}

\makeatletter
\newcommand{\specID}[1]{\def\@specID{#1}}
\newcommand{\specTitle}[1]{\def\@specTitle{#1}}
\newcommand{\specStatus}[1]{\def\@specStatus{#1}}
\newcommand{\specVersion}[1]{\def\@specVersion{#1}}
\newcommand{\specOwner}[1]{\def\@specOwner{#1}}
\newcommand{\specTraces}[1]{\def\@specTraces{#1}}

\specID{SPEC-XXX}
\specTitle{Untitled Spec}
\specStatus{DRAFT}
\specVersion{0.1}
\specOwner{Jeremie Nunez}
\specTraces{}

\newcommand{\makespecheader}{%
  \begin{center}
    {\LARGE\bfseries\color{specBlue}\@specTitle}\\[0.4em]
    {\small\color{specGray}\texttt{\@specID} \quad v\@specVersion \quad \textsc{\@specStatus} \quad \@specOwner}
  \end{center}
  \vspace{0.4em}
  \hrule
  \vspace{0.8em}
}

\pagestyle{fancy}
\fancyhf{}
\fancyhead[L]{\small\color{specGray}\@specID}
\fancyhead[R]{\small\color{specGray}\@specTitle}
\fancyfoot[C]{\small\color{specGray}\thepage}
\renewcommand{\headrulewidth}{0pt}
\makeatother

\newtcolorbox{invariant}{
  colback=specBlue!5, colframe=specBlue!75,
  title=Invariant, fonttitle=\bfseries, breakable,
  boxrule=0.5pt, arc=2pt
}

\newtcolorbox{scenario}[1]{
  colback=white, colframe=specBlue!50,
  title={Scenario: #1}, fonttitle=\bfseries, breakable,
  boxrule=0.5pt, arc=2pt
}

\newtcolorbox{ac}[1]{
  colback=specAccent!5, colframe=specAccent!60,
  title={Acceptance Criterion #1}, fonttitle=\bfseries, breakable,
  boxrule=0.5pt, arc=2pt
}

\newtcolorbox{nongoal}{
  colback=specGray!8, colframe=specGray!60,
  title=Non-Goal, fonttitle=\bfseries, breakable,
  boxrule=0.5pt, arc=2pt
}

\lstdefinestyle{codestyle}{
  basicstyle=\ttfamily\small,
  breaklines=true,
  showstringspaces=false,
  frame=single,
  framesep=4pt,
  rulecolor=\color{specBlue!30},
  backgroundcolor=\color{specBlue!2},
  xleftmargin=0.5em,
  xrightmargin=0.5em,
  keywordstyle=\color{specBlue}\bfseries,
  commentstyle=\color{specGray}\itshape,
  stringstyle=\color{specGreen}
}
\lstset{
  inputencoding=utf8,
  extendedchars=true,
  literate=
    {—}{{--}}1
    {–}{{-}}1
    {…}{{...}}1
    {’}{{'}}1
    {‘}{{'}}1
    {“}{{``}}1
    {”}{{''}}1
    {é}{{\'e}}1
    {è}{{\`e}}1
    {ê}{{\^e}}1
    {à}{{\`a}}1
    {â}{{\^a}}1
    {ç}{{\c{c}}}1
    {ô}{{\^o}}1
    {→}{{$\to$}}1
    {←}{{$\leftarrow$}}1
    {×}{{$\times$}}1
    {≥}{{$\geq$}}1
    {≤}{{$\leq$}}1
    {≠}{{$\neq$}}1
    {∈}{{$\in$}}1
    {⌈}{{$\lceil$}}1
    {⌉}{{$\rceil$}}1
    {∞}{{$\infty$}}1
    {α}{{$\alpha$}}1
    {β}{{$\beta$}}1
}
\lstdefinelanguage{TypeScript}{
  keywords={const,let,var,function,return,if,else,for,while,class,interface,type,extends,implements,import,export,from,as,async,await,new,this,true,false,null,undefined,void,any,unknown,never,string,number,boolean,readonly,private,public,protected,enum},
  sensitive=true,
  comment=[l]{//},
  morecomment=[s]{/*}{*/},
  morestring=[b]",
  morestring=[b]',
  morestring=[b]`
}
\lstset{style=codestyle}

\newcommand{\Given}{\textbf{\color{specBlue}Given} }
\newcommand{\When}{\textbf{\color{specBlue}When} }
\newcommand{\Then}{\textbf{\color{specBlue}Then} }
\providecommand{\And}{}
\renewcommand{\And}{\textbf{\color{specBlue}And} }
 (adjust depth).

\documentclass[11pt,a4paper]{article}

\usepackage[utf8]{inputenc}
\usepackage[T1]{fontenc}
\usepackage{lmodern}
\usepackage[a4paper,margin=2.3cm]{geometry}
\usepackage{parskip}
\usepackage{microtype}
\usepackage[dvipsnames]{xcolor}
\usepackage{enumitem}
\usepackage{listings}
\usepackage{tcolorbox}
\usepackage{titlesec}
\usepackage{fancyhdr}
\usepackage{array}
\usepackage{longtable}
\usepackage{booktabs}
\usepackage{amsmath}
\usepackage{amssymb}
\usepackage{hyperref}
\usepackage{cleveref}

\tcbuselibrary{skins,breakable}

\definecolor{specBlue}{RGB}{30,64,132}
\definecolor{specGray}{RGB}{90,90,90}
\definecolor{specAccent}{RGB}{198,40,40}
\definecolor{specGreen}{RGB}{30,110,60}

\hypersetup{
  colorlinks=true,
  linkcolor=specBlue,
  urlcolor=specBlue,
  citecolor=specBlue,
  pdfborder={0 0 0}
}

\titleformat{\section}{\Large\bfseries\color{specBlue}}{\thesection}{0.75em}{}
\titleformat{\subsection}{\large\bfseries\color{specBlue!80}}{\thesubsection}{0.6em}{}

\makeatletter
\newcommand{\specID}[1]{\def\@specID{#1}}
\newcommand{\specTitle}[1]{\def\@specTitle{#1}}
\newcommand{\specStatus}[1]{\def\@specStatus{#1}}
\newcommand{\specVersion}[1]{\def\@specVersion{#1}}
\newcommand{\specOwner}[1]{\def\@specOwner{#1}}
\newcommand{\specTraces}[1]{\def\@specTraces{#1}}

\specID{SPEC-XXX}
\specTitle{Untitled Spec}
\specStatus{DRAFT}
\specVersion{0.1}
\specOwner{Jeremie Nunez}
\specTraces{}

\newcommand{\makespecheader}{%
  \begin{center}
    {\LARGE\bfseries\color{specBlue}\@specTitle}\\[0.4em]
    {\small\color{specGray}\texttt{\@specID} \quad v\@specVersion \quad \textsc{\@specStatus} \quad \@specOwner}
  \end{center}
  \vspace{0.4em}
  \hrule
  \vspace{0.8em}
}

\pagestyle{fancy}
\fancyhf{}
\fancyhead[L]{\small\color{specGray}\@specID}
\fancyhead[R]{\small\color{specGray}\@specTitle}
\fancyfoot[C]{\small\color{specGray}\thepage}
\renewcommand{\headrulewidth}{0pt}
\makeatother

\newtcolorbox{invariant}{
  colback=specBlue!5, colframe=specBlue!75,
  title=Invariant, fonttitle=\bfseries, breakable,
  boxrule=0.5pt, arc=2pt
}

\newtcolorbox{scenario}[1]{
  colback=white, colframe=specBlue!50,
  title={Scenario: #1}, fonttitle=\bfseries, breakable,
  boxrule=0.5pt, arc=2pt
}

\newtcolorbox{ac}[1]{
  colback=specAccent!5, colframe=specAccent!60,
  title={Acceptance Criterion #1}, fonttitle=\bfseries, breakable,
  boxrule=0.5pt, arc=2pt
}

\newtcolorbox{nongoal}{
  colback=specGray!8, colframe=specGray!60,
  title=Non-Goal, fonttitle=\bfseries, breakable,
  boxrule=0.5pt, arc=2pt
}

\lstdefinestyle{codestyle}{
  basicstyle=\ttfamily\small,
  breaklines=true,
  showstringspaces=false,
  frame=single,
  framesep=4pt,
  rulecolor=\color{specBlue!30},
  backgroundcolor=\color{specBlue!2},
  xleftmargin=0.5em,
  xrightmargin=0.5em,
  keywordstyle=\color{specBlue}\bfseries,
  commentstyle=\color{specGray}\itshape,
  stringstyle=\color{specGreen}
}
\lstset{
  inputencoding=utf8,
  extendedchars=true,
  literate=
    {—}{{--}}1
    {–}{{-}}1
    {…}{{...}}1
    {’}{{'}}1
    {‘}{{'}}1
    {“}{{``}}1
    {”}{{''}}1
    {é}{{\'e}}1
    {è}{{\`e}}1
    {ê}{{\^e}}1
    {à}{{\`a}}1
    {â}{{\^a}}1
    {ç}{{\c{c}}}1
    {ô}{{\^o}}1
    {→}{{$\to$}}1
    {←}{{$\leftarrow$}}1
    {×}{{$\times$}}1
    {≥}{{$\geq$}}1
    {≤}{{$\leq$}}1
    {≠}{{$\neq$}}1
    {∈}{{$\in$}}1
    {⌈}{{$\lceil$}}1
    {⌉}{{$\rceil$}}1
    {∞}{{$\infty$}}1
    {α}{{$\alpha$}}1
    {β}{{$\beta$}}1
}
\lstdefinelanguage{TypeScript}{
  keywords={const,let,var,function,return,if,else,for,while,class,interface,type,extends,implements,import,export,from,as,async,await,new,this,true,false,null,undefined,void,any,unknown,never,string,number,boolean,readonly,private,public,protected,enum},
  sensitive=true,
  comment=[l]{//},
  morecomment=[s]{/*}{*/},
  morestring=[b]",
  morestring=[b]',
  morestring=[b]`
}
\lstset{style=codestyle}

\newcommand{\Given}{\textbf{\color{specBlue}Given} }
\newcommand{\When}{\textbf{\color{specBlue}When} }
\newcommand{\Then}{\textbf{\color{specBlue}Then} }
\providecommand{\And}{}
\renewcommand{\And}{\textbf{\color{specBlue}And} }
 (adjust depth).

\documentclass[11pt,a4paper]{article}

\usepackage[utf8]{inputenc}
\usepackage[T1]{fontenc}
\usepackage{lmodern}
\usepackage[a4paper,margin=2.3cm]{geometry}
\usepackage{parskip}
\usepackage{microtype}
\usepackage[dvipsnames]{xcolor}
\usepackage{enumitem}
\usepackage{listings}
\usepackage{tcolorbox}
\usepackage{titlesec}
\usepackage{fancyhdr}
\usepackage{array}
\usepackage{longtable}
\usepackage{booktabs}
\usepackage{amsmath}
\usepackage{amssymb}
\usepackage{hyperref}
\usepackage{cleveref}

\tcbuselibrary{skins,breakable}

\definecolor{specBlue}{RGB}{30,64,132}
\definecolor{specGray}{RGB}{90,90,90}
\definecolor{specAccent}{RGB}{198,40,40}
\definecolor{specGreen}{RGB}{30,110,60}

\hypersetup{
  colorlinks=true,
  linkcolor=specBlue,
  urlcolor=specBlue,
  citecolor=specBlue,
  pdfborder={0 0 0}
}

\titleformat{\section}{\Large\bfseries\color{specBlue}}{\thesection}{0.75em}{}
\titleformat{\subsection}{\large\bfseries\color{specBlue!80}}{\thesubsection}{0.6em}{}

\makeatletter
\newcommand{\specID}[1]{\def\@specID{#1}}
\newcommand{\specTitle}[1]{\def\@specTitle{#1}}
\newcommand{\specStatus}[1]{\def\@specStatus{#1}}
\newcommand{\specVersion}[1]{\def\@specVersion{#1}}
\newcommand{\specOwner}[1]{\def\@specOwner{#1}}
\newcommand{\specTraces}[1]{\def\@specTraces{#1}}

\specID{SPEC-XXX}
\specTitle{Untitled Spec}
\specStatus{DRAFT}
\specVersion{0.1}
\specOwner{Jeremie Nunez}
\specTraces{}

\newcommand{\makespecheader}{%
  \begin{center}
    {\LARGE\bfseries\color{specBlue}\@specTitle}\\[0.4em]
    {\small\color{specGray}\texttt{\@specID} \quad v\@specVersion \quad \textsc{\@specStatus} \quad \@specOwner}
  \end{center}
  \vspace{0.4em}
  \hrule
  \vspace{0.8em}
}

\pagestyle{fancy}
\fancyhf{}
\fancyhead[L]{\small\color{specGray}\@specID}
\fancyhead[R]{\small\color{specGray}\@specTitle}
\fancyfoot[C]{\small\color{specGray}\thepage}
\renewcommand{\headrulewidth}{0pt}
\makeatother

\newtcolorbox{invariant}{
  colback=specBlue!5, colframe=specBlue!75,
  title=Invariant, fonttitle=\bfseries, breakable,
  boxrule=0.5pt, arc=2pt
}

\newtcolorbox{scenario}[1]{
  colback=white, colframe=specBlue!50,
  title={Scenario: #1}, fonttitle=\bfseries, breakable,
  boxrule=0.5pt, arc=2pt
}

\newtcolorbox{ac}[1]{
  colback=specAccent!5, colframe=specAccent!60,
  title={Acceptance Criterion #1}, fonttitle=\bfseries, breakable,
  boxrule=0.5pt, arc=2pt
}

\newtcolorbox{nongoal}{
  colback=specGray!8, colframe=specGray!60,
  title=Non-Goal, fonttitle=\bfseries, breakable,
  boxrule=0.5pt, arc=2pt
}

\lstdefinestyle{codestyle}{
  basicstyle=\ttfamily\small,
  breaklines=true,
  showstringspaces=false,
  frame=single,
  framesep=4pt,
  rulecolor=\color{specBlue!30},
  backgroundcolor=\color{specBlue!2},
  xleftmargin=0.5em,
  xrightmargin=0.5em,
  keywordstyle=\color{specBlue}\bfseries,
  commentstyle=\color{specGray}\itshape,
  stringstyle=\color{specGreen}
}
\lstset{
  inputencoding=utf8,
  extendedchars=true,
  literate=
    {—}{{--}}1
    {–}{{-}}1
    {…}{{...}}1
    {’}{{'}}1
    {‘}{{'}}1
    {“}{{``}}1
    {”}{{''}}1
    {é}{{\'e}}1
    {è}{{\`e}}1
    {ê}{{\^e}}1
    {à}{{\`a}}1
    {â}{{\^a}}1
    {ç}{{\c{c}}}1
    {ô}{{\^o}}1
    {→}{{$\to$}}1
    {←}{{$\leftarrow$}}1
    {×}{{$\times$}}1
    {≥}{{$\geq$}}1
    {≤}{{$\leq$}}1
    {≠}{{$\neq$}}1
    {∈}{{$\in$}}1
    {⌈}{{$\lceil$}}1
    {⌉}{{$\rceil$}}1
    {∞}{{$\infty$}}1
    {α}{{$\alpha$}}1
    {β}{{$\beta$}}1
}
\lstdefinelanguage{TypeScript}{
  keywords={const,let,var,function,return,if,else,for,while,class,interface,type,extends,implements,import,export,from,as,async,await,new,this,true,false,null,undefined,void,any,unknown,never,string,number,boolean,readonly,private,public,protected,enum},
  sensitive=true,
  comment=[l]{//},
  morecomment=[s]{/*}{*/},
  morestring=[b]",
  morestring=[b]',
  morestring=[b]`
}
\lstset{style=codestyle}

\newcommand{\Given}{\textbf{\color{specBlue}Given} }
\newcommand{\When}{\textbf{\color{specBlue}When} }
\newcommand{\Then}{\textbf{\color{specBlue}Then} }
\providecommand{\And}{}
\renewcommand{\And}{\textbf{\color{specBlue}And} }


\specID{SPEC-TH-001}
\specTitle{Orchestrator --- query decomposition, cortex dispatch, and result assembly}
\specStatus{DRAFT}
\specVersion{0.1}
\specOwner{Jeremie Nunez}

\begin{document}
\makespecheader

\section{Context}
The Thalamus orchestrator is the top-level coordinator of the multi-cortex research agent. It receives a user query, decomposes it into a plan of cortex invocations, dispatches those invocations (in parallel across cortices, sequentially within a cortex chain), merges structured findings, and streams events to the caller. The orchestrator is deterministic glue around three moving parts: the \emph{planner} (LLM, selects cortices and steps), the \emph{cortex executor} (runs a cortex, see SPEC-TH-003), and the \emph{result assembler} (normalises, deduplicates, and streams findings).

\section{Goals}
\begin{itemize}[leftmargin=1.2em]
  \item Turn a single query into a plan of N missions (one per selected cortex) with well-formed steps.
  \item Execute missions in parallel while preserving per-mission step order and chaining (\texttt{pickFrom} / \texttt{feedsInto}).
  \item Emit a typed event stream (\texttt{tool\_start}, \texttt{tool\_result}, \texttt{tool\_error}, \texttt{agent\_done}) that the UI or upstream caller can consume incrementally.
  \item Enforce mutation guards: tools that mutate state only run when the query intent explicitly authorises mutation.
  \item Enforce per-tool timeouts and make a failed step non-fatal for the remainder of its mission.
\end{itemize}

\section{Non-Goals}
\begin{nongoal}
The orchestrator does not implement domain logic (scoring, SQL, LLM reasoning) --- that is owned by each cortex (SPEC-TH-003). It does not persist findings to the knowledge graph; the caller is responsible for consuming the event stream and committing. It does not negotiate budget targets across iterations; that loop lives one level above.
\end{nongoal}

\section{Invariants}
\begin{invariant}
For every mission \texttt{M} and every step \texttt{S} in \texttt{M}, exactly one of \texttt{tool\_result} or \texttt{tool\_error} is emitted, and it references a \texttt{toolCallId} previously seen in a \texttt{tool\_start}. Exactly one \texttt{agent\_done} is emitted per mission, after all its step events.
\end{invariant}

\begin{invariant}
A mutating tool is dispatched if and only if the caller's intent is in the set of action intents. A blocked mutation yields a \texttt{tool\_error} with reason \texttt{mutation\_blocked} and does not call the tool.
\end{invariant}

\begin{invariant}
A step timeout or thrown exception in one step never prevents subsequent steps in the same mission from being attempted, and never cancels sibling missions.
\end{invariant}

\section{Interface}
Public surface consumed by the HTTP layer and the recursive research loop.

\begin{lstlisting}[language=TypeScript]
export interface Mission {
  agent: AgentRole;
  steps: MissionStep[];
}

export interface MissionStep {
  tool: string;
  params: Record<string, unknown>;
  pickFrom?: PickStrategy;      // select from previous result array
  feedsInto?: string;           // param name fed by previous id(s)
  size?: CardSize;              // UI hint only
}

export type PipelineSSEEvent =
  | ToolStartEvent
  | ToolResultEvent
  | ToolErrorEvent
  | AgentDoneEvent;

export class MissionExecutor {
  constructor(toolMap: Map<string, DynamicStructuredTool>);

  executeAllMissions(
    missions: Mission[],
    intent: string,
  ): AsyncGenerator<PipelineSSEEvent>;

  executeMission(
    mission: Mission,
    intent: string,
  ): AsyncGenerator<PipelineSSEEvent>;

  retrySingleTool(
    toolName: string,
    modifiedInput: Record<string, unknown>,
    agent?: AgentRole,
  ): AsyncGenerator<PipelineSSEEvent>;
}
\end{lstlisting}

\section{Scenarios}

\begin{scenario}{Nominal parallel dispatch}
\Given a plan of three missions \texttt{[M1, M2, M3]} and a read-only intent \\
\When \texttt{executeAllMissions} is consumed to completion \\
\Then step events from the three missions may interleave in time order \\
\And each mission emits exactly one \texttt{agent\_done} after its last step \\
\And the total number of \texttt{tool\_start} events equals the total number of steps.
\end{scenario}

\begin{scenario}{Step chaining via pickFrom/feedsInto}
\Given a mission whose step 1 returns \verb|{ items: [{id:7}, {id:9}] }| \\
\And step 2 declares \texttt{pickFrom: "top\_by\_relevance"} and \texttt{feedsInto: "itemId"} \\
\When the mission is executed \\
\Then step 2 is invoked with \texttt{params.itemId === 7} \\
\And no \texttt{tool\_error} is emitted for step 2 due to chaining.
\end{scenario}

\begin{scenario}{Mutation blocked on read intent}
\Given a mission whose step invokes a mutating tool \\
\And the caller intent is not in the action intents set \\
\When the mission is executed \\
\Then a single \texttt{tool\_error} with reason \texttt{mutation\_blocked} is emitted for that step \\
\And the tool is never invoked \\
\And subsequent steps in the mission still run.
\end{scenario}

\begin{scenario}{Tool timeout is recoverable}
\Given a mission with three steps where step 2 exceeds its timeout \\
\When the mission is executed \\
\Then step 2 emits \texttt{tool\_error} with a timeout message \\
\And step 3 still emits \texttt{tool\_start} and terminates normally \\
\And sibling missions are unaffected.
\end{scenario}

\begin{scenario}{Empty upstream breaks the chain}
\Given step 1 returns an empty result array and step 2 declares \texttt{pickFrom} \\
\When the mission is executed \\
\Then step 2 emits \texttt{tool\_error} with reason \texttt{no\_results\_from\_previous\_step} \\
\And no further steps of that mission are attempted.
\end{scenario}

\section{Acceptance Criteria}

\begin{ac}{AC-1}
Given a plan of $n$ missions, \texttt{executeAllMissions} yields exactly $n$ \texttt{agent\_done} events; for each mission the \texttt{agent\_done} is the last event bearing that \texttt{agent} field.
\end{ac}

\begin{ac}{AC-2}
For every step, the \texttt{toolCallId} appearing in its \texttt{tool\_start} matches the \texttt{toolCallId} of its terminal \texttt{tool\_result} or \texttt{tool\_error}.
\end{ac}

\begin{ac}{AC-3}
A step marked as using a mutating tool, under a non-action intent, produces a \texttt{tool\_error} with code \texttt{mutation\_blocked} and does not invoke the tool (verified via a spy).
\end{ac}

\begin{ac}{AC-4}
A step whose underlying tool rejects or exceeds the applicable timeout emits exactly one \texttt{tool\_error} and does not prevent later steps of the same mission from emitting.
\end{ac}

\begin{ac}{AC-5}
Chaining with \texttt{pickFrom: "top\_by\_relevance"} feeds \texttt{results[0].id} into the declared \texttt{feedsInto} parameter of the next step.
\end{ac}

\begin{ac}{AC-6}
An exception thrown by a tool in mission $M_i$ does not cancel or short-circuit mission $M_j$ for $j \neq i$: both missions still emit an \texttt{agent\_done}.
\end{ac}

\begin{ac}{AC-7}
\texttt{retrySingleTool} invoked with a tool in the retry-blocked set emits a single \texttt{tool\_error} and does not invoke the tool.
\end{ac}

\section{Traceability}

\begin{center}
\begin{tabular}{@{}lll@{}}
\toprule
AC & Test file & Test name \\
\midrule
AC-1 & \texttt{packages/thalamus/tests/orchestrator.spec.ts} & \texttt{emits one agent\_done per mission} \\
AC-2 & \texttt{packages/thalamus/tests/orchestrator.spec.ts} & \texttt{toolCallId round-trip} \\
AC-3 & \texttt{packages/thalamus/tests/orchestrator.spec.ts} & \texttt{mutation blocked on read intent} \\
AC-4 & \texttt{packages/thalamus/tests/orchestrator.spec.ts} & \texttt{timeout does not abort mission} \\
AC-5 & \texttt{packages/thalamus/tests/orchestrator.spec.ts} & \texttt{pickFrom feeds top result id} \\
AC-6 & \texttt{packages/thalamus/tests/orchestrator.spec.ts} & \texttt{sibling mission isolation} \\
AC-7 & \texttt{packages/thalamus/tests/orchestrator.spec.ts} & \texttt{retry blocks disallowed tools} \\
\bottomrule
\end{tabular}
\end{center}

\section{Open Questions}
\begin{itemize}[leftmargin=1.2em]
  \item Should \texttt{agent\_done} carry per-mission aggregate metrics (step count, error count) or remain a pure marker?
  \item Should chaining support multi-source fan-in (step $k$ reads from step $k-2$), or is single-predecessor sufficient for v1?
  \item Where is the hard cap on concurrent missions enforced --- orchestrator, caller, or transport layer?
\end{itemize}

\end{document}
